\section{Methodology and Experimental Setup}\label{methodology-and-experimental-setup}

\subsection{Checkpoint Audit and Method Starting Point}\label{checkpoint-audit-and-method-starting-point}

All experiments are implemented in Python with PyTorch on top of the CellSAM codebase \cite{israel2024cellsam}. Training is conducted on the ALICE High-Performance Computing (HPC) cluster at Leiden University, mainly on NVIDIA L4 GPUs (24 GB VRAM). Local machines are used for result auditing, visualization, and baseline reruns.

Before defining the thesis mainline, we audited the released CellSAM checkpoint as a software artifact. The public checkpoint contains three relevant branches: \texttt{model}, \texttt{model\_cp}, and \texttt{cellfinder}. For segmentation in this thesis, we consistently use the released \texttt{model\_cp} single-mask route (\texttt{multimask\_output=False}) as the starting point and as the inference branch.

The audit also defines a clear interpretation boundary. Parameter-level comparisons show that \texttt{model} and \texttt{model\_cp} are globally different branches, so branch identity must be explicit whenever results are reported. We therefore keep only verifiable implementation facts in the main text and move full tensor-level comparisons to Appendix~A.

It is also important to distinguish architecture construction from weight loading. In the released code, \texttt{self.model\ =\ sam\_model\_registry{[}"vit\_b"{]}()} is instantiated first and \texttt{self.model\_cp} is then created as a deep copy of this freshly constructed model. Thus, before the released CellSAM checkpoint is loaded, the two branches share the SAM ViT-B architecture and initialization state \cite{kirillov2023sam}, but they are not automatically the original Meta pretrained SAM checkpoint.

\subsection{T28 Mainline Training Setup}\label{t28-mainline-training-setup}

The current thesis segmentation mainline follows the T28 configuration. Instead of treating the released CellSAM checkpoint as a simple neck-only continuation of Stage 1, we explicitly use the released \texttt{model\_cp} segmentation branch as the starting point for fine-tuning. In the main configuration, the image encoder and prompt encoder are frozen, and only the mask decoder is updated. Input images are resized to \texttt{1024\ x\ 1024}, and a cardiomyocyte-aware three-channel input is used (\texttt{[BF,\ Actn2,\ DAPI]}), i.e., the legacy T28 channel order retained for sealed-route consistency.

The training objective in the current mainline is not the older ``Best Config'' loss. It is a \texttt{CombinedLoss} with Dice \cite{milletari2016vnet}, BCE, Boundary \cite{kervadec2019boundary}, AJI \cite{kumar2017aji}, and Focal \cite{lin2017focal} components, plus an explicit IoU-head regression term added in the training loop (aligned with SAM mask-quality prediction semantics \cite{kirillov2023sam}). This distinction matters because the thesis now treats T28, rather than Best Config, as the primary segmentation route.

Table 4.1 summarizes the main training configuration used for the sealed T28 route.

\textbf{Table 4.1: Main training configuration used in T28}

{%
\begin{longtable}[]{@{}ll@{}}
\toprule\noalign{}
Hyperparameter & Value \\
\midrule\noalign{}
\endhead
\bottomrule\noalign{}
\endlastfoot
Base checkpoint & CellSAM released checkpoint, \texttt{model\_cp} branch \\
Trainable module & SAM mask decoder \\
Frozen modules & Image encoder, prompt encoder \\
Input encoding & Three-channel cardiomyocyte-aware route (legacy T28 order), \texttt{{[}BF,\ Actn2,\ DAPI{]}} \\
Image size & \texttt{1024\ x\ 1024} \\
Batch size & 4 \\
Optimizer & AdamW \\
Learning rate & \texttt{1e-4} \\
Weight decay & \texttt{1e-4} \\
LR schedule & Cosine warmup, 5 warmup epochs \\
Max epochs & 80 \\
Early stopping & PQ-based, patience = 15 \\
Training paradigm & Instance-level, one GT cell mask per box \\
Box expansion & 10 percent \\
Main loss & \texttt{CombinedLoss} \\
Extra loss & IoU-head MSE, weight = 0.1 \\
Active loss terms & Dice + BCE, Boundary, AJI, Focal \\
Disabled loss terms & Contour, Topology, Size \\
Mixed precision & AMP \\
\end{longtable}
}

For completeness, the active T28 loss configuration is:

\begin{itemize}
\tightlist
\item
  \texttt{pos\_weight\ =\ 10}
\item
  \texttt{boundary\_weight\ =\ 0.3}
\item
  \texttt{aji\_weight\ =\ 0.2}
\item
  \texttt{focal\_weight\ =\ 0.3}
\item
  \texttt{contour\_weight\ =\ 0.0}
\item
  \texttt{iou\_weight\ =\ 0.1}
\end{itemize}

This is the configuration that should be described in the thesis method chapter. Earlier settings such as Phase 1 or Best Config remain useful as historical comparisons, but they are no longer the default writing reference.

\subsection{Final Inference Pipeline}\label{final-inference-pipeline}

The finalized deployment-oriented pipeline used in this thesis is:

\begin{enumerate}
\def\labelenumi{\arabic{enumi}.}
\tightlist
\item
  Build the cardiomyocyte-aware input tensor and resize to \texttt{1024\ x\ 1024}.
\item
  Generate prompt boxes from the selected detector route (\texttt{T33g}, \texttt{dapi\_cm} candidate-aware mode) for end-to-end evaluation, or use GT boxes for Oracle evaluation.
\item
  Run per-box single-mask segmentation using the T28-adapted \texttt{model\_cp} branch.
\item
  Resolve multi-instance overlap conflicts with \texttt{argmax\_prob} assignment.
\item
  Apply instance-level postprocessing and assemble final instance masks.
\item
  Compute split-aware metrics and report Oracle versus end-to-end results separately.
\end{enumerate}

Table 4.2 fixes the module-role mapping to avoid ambiguity between detector and segmenter responsibilities.

\textbf{Table 4.2: Final module-role mapping in the thesis pipeline}

{%
\begin{longtable}[]{@{}lll@{}}
\toprule\noalign{}
Stage & Module & Role \\
\midrule\noalign{}
\endhead
\bottomrule\noalign{}
\endlastfoot
Input preparation & Semantic channel mapping + resizing & Construct cardiomyocyte-aware model input at \texttt{1024\ x\ 1024} \\
Prompt generation & \texttt{T33g} detector (\texttt{dapi\_cm}, candidate-aware) & Provide automatic box prompts for deployment \\
Segmentation core & \texttt{T28} on released \texttt{model\_cp} branch & Predict per-box cell masks (\texttt{multimask\_output=False}) \\
Conflict handling & \texttt{argmax\_prob} assignment & Resolve overlapping predictions at pixel level \\
Postprocessing & Morphology + largest connected component & Smooth masks and remove implausible fragments \\
Evaluation & Unified instance-metric protocol & Report PQ/BM-Dice/AJI plus split-aware F1/RQ evidence \\
\end{longtable}
}

\subsection{Preprocessing and Postprocessing Design}\label{preprocessing-and-postprocessing-design}

\textbf{Preprocessing.}
The segmentation input is standardized to \texttt{1024\ x\ 1024}. For the sealed T28 route, channels follow the cardiomyocyte-aware three-channel setting \texttt{[BF,\ Actn2,\ DAPI]} with per-channel normalization in the data pipeline. Oracle and end-to-end branches share the same downstream segmentation preprocessing, while they differ only in prompt source (GT boxes versus detector boxes).

\textbf{Prompt geometry handling.}
During training, each GT box is expanded by \texttt{10\%} before loss computation to reduce sensitivity to tight-box truncation. During inference, box clipping with the same expansion ratio is enabled by default in the thesis mainline and applied before conflict resolution; the unclipped setting is used only as a supplementary comparison condition. This keeps geometric support consistent between training and deployment-time mask aggregation.

\textbf{Postprocessing.}
Final instance masks are postprocessed per instance using morphology operations (small-hole/island cleanup, opening/closing, dilation-erosion smoothing, Gaussian boundary smoothing), followed by largest-connected-component retention. Optional size validation is available (\texttt{min\_cell\_area=13884}, \texttt{max\_cell\_area=174735} at \texttt{1024\ x\ 1024}) but is disabled in the default thesis evaluation route.

\subsection{Optimization Objectives}\label{optimization-objectives}

The T28 training objective combines segmentation and IoU-head regression:
\begin{equation}
\mathcal{L}_{\mathrm{total}} = \mathcal{L}_{\mathrm{combined}} + \lambda_{\mathrm{iou}}\mathcal{L}_{\mathrm{iou}}, \qquad \lambda_{\mathrm{iou}} = 0.1.
\end{equation}

In implementation, the segmentation base term is:
\begin{equation}
\mathcal{L}_{\mathrm{base}} = 0.5\mathcal{L}_{\mathrm{Dice}} + 0.5\mathcal{L}_{\mathrm{BCE}},
\end{equation}
and the active auxiliary family for T28 is
\begin{equation}
\mathcal{K}=\{\mathrm{Boundary},\ \mathrm{AJI},\ \mathrm{Focal}\}.
\end{equation}
With normalized weighting,
\begin{equation}
\mathcal{L}_{\mathrm{combined}} =
\frac{w_{\mathrm{base}}}{T}\mathcal{L}_{\mathrm{base}} +
\sum_{k\in\mathcal{K}}\frac{w_k}{T}\mathcal{L}_k, \qquad
T = w_{\mathrm{base}} + \sum_{k\in\mathcal{K}} w_k.
\end{equation}
The active terms correspond to Dice \cite{milletari2016vnet}, Boundary \cite{kervadec2019boundary}, AJI-inspired optimization \cite{kumar2017aji}, and Focal loss \cite{lin2017focal}.

For detector training (T33 family), matching is solved by Hungarian assignment \cite{kuhn1955hungarian} and the objective follows the AnchorDETR-style weighted combination used in the CellSAM codebase \cite{israel2024cellsam}:
\begin{equation}
\mathcal{L}_{\mathrm{det}} = 2\mathcal{L}_{\mathrm{ce}} + 5\mathcal{L}_{\mathrm{bbox}} + 2\mathcal{L}_{\mathrm{giou}},
\end{equation}
with auxiliary decoder-layer losses at the same weight pattern.

\subsection{Operational Algorithms}\label{operational-algorithms}

To make the methodology executable and reproducible at implementation level, we explicitly provide two algorithms in the main text: one for T28 decoder-focused training under Oracle prompts, and one for end-to-end inference with automatic prompts.

\begin{algorithm}[t]
\caption{T28 Decoder-Focused Training and Oracle Evaluation}
\label{alg:t28_train_oracle}
\begin{algorithmic}[1]
\Require Training set $\mathcal{D}_{train}=\{(I,M_{gt},B_{gt})\}$, validation set $\mathcal{D}_{val}$, released checkpoint $C$
\Ensure Best checkpoint $W^\star$ and Oracle validation history $\mathcal{H}_{val}$
\State Load $C$ and select segmentation branch \texttt{model\_cp}
\State Freeze image encoder and prompt encoder; set mask decoder trainable
\For{epoch $e=1$ to $E_{max}$}
  \For{mini-batch $(I,M_{gt},B_{gt}) \in \mathcal{D}_{train}$}
    \State Build three-channel cardiomyocyte-aware input (\texttt{[BF, Actn2, DAPI]})
    \State Use $B_{gt}$ as Oracle prompts and predict masks with single-mask path
    \State Compute $L_{main}$ (Dice+BCE, Boundary, AJI, Focal)
    \State Compute $L_{iou}$ (IoU-head MSE) and $L \gets L_{main} + \lambda_{iou} L_{iou}$
    \State Update mask decoder parameters with AdamW
  \EndFor
  \State Evaluate on $\mathcal{D}_{val}$ under Oracle prompts and record metrics
  \If{validation PQ improves}
    \State Save current checkpoint as best
  \EndIf
  \If{early stopping criterion is met}
    \State \textbf{break}
  \EndIf
\EndFor
\State \Return $W^\star, \mathcal{H}_{val}$
\end{algorithmic}
\end{algorithm}

\begin{algorithm}[t]
\caption{End-to-End Cardiomyocyte Instance Segmentation with Automatic Prompts}
\label{alg:e2e_prompt_seg}
\begin{algorithmic}[1]
\Require Image $I$ (BF, Actn2, DAPI), trained checkpoint $W^\star$, prompt generator $G$, inference config $\Psi$
\Ensure Final instance mask $M_{final}$ and optional end-to-end metrics
\State Load T28 checkpoint $W^\star$ on \texttt{model\_cp} single-mask path
\State Generate automatic prompt boxes: $B_{auto} \gets G(I)$
\State Build three-channel cardiomyocyte-aware input from $I$
\For{each prompt box $b \in B_{auto}$}
  \State Predict binary mask $m_b$ with T28
\EndFor
\State Aggregate predicted masks $\{m_b\}$ and apply postprocess pipeline $\Psi$
\State Obtain final instance mask $M_{final}$
\If{ground truth is available}
  \State Compute end-to-end metrics (PQ, BM-Dice, AJI, SQ, RQ/F1)
\EndIf
\State \Return $M_{final}$ and metrics
\end{algorithmic}
\end{algorithm}

\subsection{Evaluation Protocol}\label{evaluation-protocol}

We use two evaluation settings.

\textbf{Oracle evaluation.} Ground-truth bounding boxes are used as prompts. This isolates segmentation quality from detection errors and is the main evaluation mode for prompt-based CellSAM variants such as CellSAM-pretrained, SAM ViT-B, MedSAM, and the thesis mainline T28.

\textbf{End-to-end evaluation.} Automatically detected bounding boxes are used as prompts. In the current project, this setting is represented by DAPI/Adaptive detector routes and the selected \texttt{T33g(dapi\_cm,\ candidate\_aligned,\ q35)+T28} line. The gap between Oracle and end-to-end performance quantifies how much the detection stage limits deployment quality.

All final tables must explicitly label the data split. In particular, \texttt{val71} and \texttt{test73} results must not be mixed in the same table without a split column or a separate subsection.

The sealed T28 mainline result is reported as dual-seed mean on \texttt{val71}. Locked \texttt{test73} single-checkpoint numbers (for example T27a) are retained as reference tier and are never treated as direct substitutes for the T28 mainline summary. Likewise, preliminary single-seed experiments are not treated as final main-table evidence.

\subsubsection{Evaluation Metrics}\label{evaluation-metrics}

We report the following metrics.

\begin{itemize}
\tightlist
\item
  \textbf{PQ@0.5}: Panoptic Quality, defined as \texttt{PQ\ =\ SQ\ x\ RQ}. This is the primary metric because it jointly reflects segmentation quality and recognition quality.
\item
  \textbf{BM-1to1 Dice}: Hungarian one-to-one best-match Dice between predicted and ground-truth instances.
\item
  \textbf{AJI}: Aggregated Jaccard Index, used as an instance-level overlap measure.
\item
  \textbf{SQ}: Segmentation Quality of matched pairs.
\item
  \textbf{RQ}: Recognition Quality.
\end{itemize}

When the same TP/FP/FN definition is used, \texttt{RQ} is mathematically equivalent to \texttt{F1}. However, some project scripts output \texttt{per-image\ mean\ RQ} and \texttt{global\ micro-F1} separately. Therefore, when both are reported, the aggregation scheme must be stated explicitly.

\subsection{Baseline Methods and Reporting Scope}\label{baseline-methods-and-reporting-scope}

The thesis main comparison currently uses only locked baselines:

\begin{itemize}
\tightlist
\item
  \textbf{Cellpose (T31)}: rerun with the latest Cellpose v4.0.1 baseline setting, using \texttt{diameter\ =\ 250}.
\item
  \textbf{SAM ViT-B}: original SAM without cell-specific fine-tuning, evaluated with GT box prompts.
\item
  \textbf{CellSAM (pretrained)}: released CellSAM checkpoint without cardiomyocyte fine-tuning, evaluated on the released \texttt{model\_cp} segmentation branch.
\item
  \textbf{MedSAM}: pretrained medical SAM baseline with GT box prompts.
\item
  \textbf{T28}: sealed segmentation mainline for downstream detector-driven integration (dual-seed \texttt{val71} reporting).
\end{itemize}

Historical or partially audited results are not merged into the same evidence tier:

\begin{itemize}
\tightlist
\item
  \textbf{Preliminary single-seed multi-channel runs} remain appendix-only observations unless multi-seed stability is demonstrated.
\item
  \textbf{T11/T30} are encoder-adaptation explorations and are not part of the current main result.
\item
  Older deprecated baseline routes, especially the previous brightfield-only Cellpose path, should not be cited in the thesis.
\end{itemize}

\subsection{Main Figure Briefs Locked for Writing}\label{main-figure-briefs-locked-for-writing}

To avoid drift between communication notes and the thesis body, we lock the two core method-figure briefs here before final artwork is inserted.

\textbf{Figure 2 (Method Overview, planned placement in Chapter 4).}
This figure is defined as a high-level five-stage flow: input handling, prompt source, released-checkpoint branch selection, T28 adaptation, and output/evaluation split. Mandatory elements are cardiomyocyte-aware three-channel input, Oracle vs automatic prompts, released segmentation branch (\texttt{model\_cp}), and split-aware outputs. To avoid duplication, detector internals and full training details are not expanded in the main trunk; training specifics are shown only as side notes.

\textbf{Locked caption draft for Figure 2.}
\emph{Overview of the thesis pipeline. Brightfield images are converted into a CellSAM-compatible input representation, while DAPI and Actn2 provide biological priors for prompting and detector design. Prompts are supplied either from ground-truth boxes under Oracle evaluation or from automatic detectors under end-to-end evaluation. The released CellSAM checkpoint is first audited to identify the \texttt{model\_cp} segmentation branch, after which the T28 mainline adapts the mask decoder while keeping the image encoder and prompt encoder frozen.}

\textbf{Figure 3 (Checkpoint Audit Schema, planned placement in supplementary material).}
This figure is defined as a branch-relationship audit: released checkpoint with \texttt{model}, \texttt{model\_cp}, and \texttt{cellfinder}; explicit arrow for released segmentation inference through \texttt{model\_cp}; and explicit note that \texttt{cellfinder} backbone aligns with \texttt{model.image\_encoder} non-neck body rather than \texttt{model\_cp}.

\textbf{Locked caption draft for Figure 3.}
\emph{Conceptual audit of the released CellSAM checkpoint. The public release contains separate model, model\_cp, and cellfinder branches. In the released inference implementation, segmentation is performed through the model\_cp branch, while the CellFinder backbone aligns with the non-neck body of model.image\_encoder rather than with model\_cp. The figure summarizes only implementation facts that are directly verifiable from the public release and does not claim full recovery of hidden training history.}

Final editable artwork is tracked as pending assets and will be inserted without changing these locked semantics.
